% Generated draft: requires mathematical and source review.
\documentclass[11pt]{article}
\usepackage[margin=25mm]{geometry}
\usepackage{fontspec}
\setmainfont{Latin Modern Roman}
\usepackage{amsmath,amssymb,mathtools,unicode-math}
\setmathfont{Latin Modern Math}
\usepackage{enumitem,needspace,xurl,hyperref}
\setlength{\parindent}{0pt}
\setlength{\parskip}{5pt}
\setlength{\emergencystretch}{3em}
\newcommand{\N}{\mathbb{N}}
\newcommand{\Z}{\mathbb{Z}}
\newcommand{\Q}{\mathbb{Q}}
\newcommand{\R}{\mathbb{R}}
\newcommand{\C}{\mathbb{C}}
\newcommand{\F}{\mathbb{F}}
\newcommand{\A}{\mathbb{A}}
\newcommand{\PP}{\mathbb{P}}
\newcommand{\E}{\mathbb{E}}
\newcommand{\eps}{\varepsilon}
\newcommand{\dd}{\,\mathrm{d}}
\begin{document}
\section*{UnsolvedMath ID: 800581}
\textbf{Source problem code:} catalogue entry
\subsection{Definitions and mathematical statement}
The supplied question is: ``Is Borel's Conjecture for $\mathsf{S}_{\mathrm{fin}}(\mathrm{B}_\Omega,\mathrm{B}_\Omega)$ consistent?'' The record does not define these terms or specify the underlying foundational theory. The following makes explicit the conventional sets-of-reals interpretation, which requires confirmation against the source.

Let $X\subseteq\R$, with its subspace topology. A countable Borel $\omega$-cover of $X$ is an at most countable family $\mathcal U$ of Borel subsets of $X$ such that
\[
X\notin\mathcal U
\quad\text{and}\quad
(\forall F\subseteq X\text{ finite})(\exists U\in\mathcal U)\;F\subseteq U.
\]
Write $\mathrm B_\Omega(X)$ for the family of all such covers. Say that $X$ satisfies $\mathsf S_{\mathrm{fin}}(\mathrm B_\Omega,\mathrm B_\Omega)$ if, for every sequence $(\mathcal U_n)_{n\in\N}$ with $\mathcal U_n\in\mathrm B_\Omega(X)$, there are finite subfamilies $\mathcal F_n\subseteq\mathcal U_n$ such that
\[
\bigcup_{n\in\N}\mathcal F_n\in\mathrm B_\Omega(X).
\]
Under the conventional interpretation of Borel's Conjecture for this property, the proposed assertion is
\[
\forall X\subseteq\R,\qquad
\mathsf S_{\mathrm{fin}}(\mathrm B_\Omega,\mathrm B_\Omega)(X)
\ \Longrightarrow\ X\text{ is at most countable}.
\]
The problem asks whether this assertion can hold in a model of $\mathrm{ZFC}$. In the usual relative-consistency formulation, the question is whether consistency of $\mathrm{ZFC}$ implies consistency of $\mathrm{ZFC}$ together with the displayed assertion. The ambient class of spaces, the cover conventions, the meaning of the generalized Borel conjecture, and the consistency base are not explicitly supplied in the record; this formulation is therefore conditional on those conventional interpretations.
\subsection{Short English statement}
Under the conventional interpretation, can it consistently happen that every set of reals satisfying $\mathsf S_{\mathrm{fin}}(\mathrm B_\Omega,\mathrm B_\Omega)$ is countable? This selection property requires choosing finitely many members from each of any sequence of countable Borel $\omega$-covers so that all chosen members together form an $\omega$-cover. The intended conventions require source confirmation.
\subsection{Sources}
\begin{itemize}
\item Local UnsolvedMath registry: supplied record for permanent ID $800581$, exact source code \texttt{catalogue entry}, in the AMR Open Problem Lists collection; upstream revision \texttt{d371ee917ac9647fbb06cd6fc529deec6f3dda82}. The supplied status is \texttt{open}; the record was updated on 2026-10-16. This is a supplied registry label, not an independently verified literature assessment.
\item Tsaban, \emph{Selection Principles and Special Sets of Reals: Open Problems} (2006), Problem 7. Supplied source URL: \url{https://arxiv.org/abs/math/0609601}. The registry lists Tsaban as proposer and 2006 as the proposed year.
\item The supplied attribution/context marker is \texttt{[CBC]}; its expansion and attribution details are not included in the record.
\item The supplied provenance reports access on 2026-10-16 and extraction by \texttt{arxiv\_latex}. This access information is supplied metadata; the source has not been independently visited or verified for this draft.
\end{itemize}
\end{document}
